\documentclass[12pt]{article}
\usepackage[utf8]{inputenc}
\usepackage[T1]{fontenc}
\usepackage{hyperref}
\title{Tai Laing Language Resources for K-12 Curriculum Development}
\author{Compiled by Claude Sonnet 5 (claude-sonnet-5) for Andr\'as Kornai}
\date{}
\begin{document}
\maketitle
\section{Introduction}
Tai Laing (ISO 639-3 tjl; also Tai Leng, Tai Lai) is a Tai language of Myanmar. It is spoken in Homalin Township (Sagaing Region) and in Kachin State between Bhamo and Myitkyina. One source gives about 100,000 speakers. It is closely related to Khamti and Shan and is written in a variant of the Burmese script. It is not taught in schools, but a small cultural revival is reported.
\section{Primary Educational Resources}
None found in the search results.
\section{Linguistic and Reference Materials}
\begin{itemize}
\item Wikipedia article on Tai Laing: \url{https://en.wikipedia.org/wiki/Tai_Laing_language}
\item Omniglot page on the Tai Laing script and language: \url{https://www.omniglot.com/writing/tailaing.htm}
\item Omniglot page on Tai Laing numbers: \url{https://omniglot.com/language/numbers/tailaing.htm}
\end{itemize}
\section{Religious and Community Resources}
\begin{itemize}
\item Scripture Earth listing for tjl (content not checked): \url{https://www.scriptureearth.org/00eng.php?iso=tjl}
\item Joshua Project profile: \url{https://www.joshuaproject.net/index.php/people_groups/20584}
\item Asia Harvest profile of the Tai Laing (PDF): \url{https://operation.asiaharvest.org/MYANMAR-Tai-Laing.pdf}
\end{itemize}
\section{Digital and Audio Resources}
\begin{itemize}
\item Wikitongues page for tjl (EGIDS 6b, threatened): \url{https://wikitongues.org/languages/tjl/}
\end{itemize}
\section{Curriculum Development Context}
Only an established script and general descriptions were found. No textbooks, dictionaries or school programmes turned up, so a curriculum would have to be built with the community. The dialect names (Tai Lai, Tai Nai, Tai Nine, Tai Leng) are used inconsistently across sources.
\end{document}